\section{Gradient and Directional Derivatives}

\subsection{The gradient}
For $f : \R^n \to \R$,
\[ \grad f = \left\langle \frac{\partial f}{\partial x_1}, \ldots, \frac{\partial f}{\partial x_n} \right\rangle. \]

\subsection{Directional derivative}
The rate of change of $f$ at $\vc{x}_0$ in unit direction $\vc{u}$:
\[ D_{\vc u} f(\vc x_0) = \lim_{h \to 0} \frac{f(\vc x_0 + h \vc u) - f(\vc x_0)}{h}
= \grad f(\vc x_0) \dotp \vc u \qquad \text{($f$ differentiable)}. \]
\textbf{$\vc u$ must be a unit vector.}

\subsection{Geometric meaning}
\begin{keybox}[The three properties of $\grad f$]
\begin{enumerate}[leftmargin=*,nosep]
  \item $\grad f(\vc x_0)$ points in the direction of \emph{steepest increase} of $f$ at $\vc x_0$.
  \item $\lVert \grad f(\vc x_0) \rVert$ is the maximum value of $D_{\vc u} f(\vc x_0)$.
  \item $\grad f(\vc x_0)$ is normal to the level set $\{ f = f(\vc x_0) \}$ at $\vc x_0$.
\end{enumerate}
\end{keybox}

\begin{example}
$f(x,y,z) = x^2 + y^2 - z$ at $(1,2,3)$: $\grad f = \langle 2x, 2y, -1 \rangle = \langle 2,4,-1\rangle$.
Steepest ascent direction. The tangent plane to the level surface $f = 2$ at $(1,2,3)$ is
$2(x-1) + 4(y-2) - (z-3) = 0$.
\end{example}

\subsection{Tangent plane to a level surface}
If $F(x,y,z) = c$ and $\grad F(\vc x_0) \neq 0$, the tangent plane at $\vc x_0$ is
\[ \grad F(\vc x_0) \dotp (\vc r - \vc x_0) = 0. \]

\subsection{Gradient identities}
\[
\grad (f + g) = \grad f + \grad g, \quad
\grad (f g) = f \grad g + g \grad f, \quad
\grad (f/g) = \frac{g \grad f - f \grad g}{g^2}.
\]
For a composition $h(\vc x) = \varphi(f(\vc x))$: $\grad h = \varphi'(f) \grad f$.
